\chapter*{General introduction}
\addcontentsline{toc}{chapter}{General introduction}

The aviation industry is categorised among the sectors that require reliability so much, because if the avionics system (the electronic pieces used in aircraft) fails unexpectedly, it will cause operational disruptions that demand costly maintenance interventions, and in critical cases, it can affect safety.

Modern aircraft continuously generate large volumes of data across multiple systems, such as electronic, thermal and communication subsystems.

There are three types of maintenance. The first is reactive maintenance, where maintenance is required after disruptions occur. This is not perfect anymore for the aviation field. The second type is preventive maintenance, which consists of controlling what must be controlled under a strict timetable. This is not always efficient because it can lead to unnecessary part replacements and that is a loss of money. The third type is predictive maintenance, which is based on past data. We forecast disruptions before catastrophic failure occurs.

Our solution touches this challenge. Leveraging the massive quantity of operational time-series data generated requires intelligent and automatic software solutions. So we propose the conception and development of an intelligent avionics health monitoring and predictive maintenance system. This system helps users detect anomalies during data streams, using different methods such as threshold-based detection and Isolation Forest (an AI model that will operate on generated data).

This report is organised as follows. Chapter 1 presents the company presentation. Chapter 2 presents the state of the art and project positioning. Chapter 3 describes the requirements analysis and Scrum methodology. Chapter 4 covers the design and architecture. Chapter 5 presents Sprint 1: Telemetry Simulation. Chapter 6 presents Sprint 2: Anomaly Detection Engine. Chapter 7 presents Sprint 3: Dashboard and Integration. Chapter 8 discusses testing, validation, and results. Finally, the general conclusion summarises the work and outlines future perspectives.
